% =============================================================================
% Chapter 16 -- Discussion
% =============================================================================

\chapter{Discussion}
\label{ch:discussion}

\section{Answering the Central Question}
\label{sec:disc:central}

The central research question of Part~V asks: \emph{Is class imbalance the true cause of poor classifier performance?}
Our controlled experiment yields a nuanced but affirmative answer.

\textbf{Yes, class imbalance causes performance degradation} --- but not uniformly, not linearly, and not irreversibly.
The degradation is real and consistent across all datasets and metrics once $\IR > 1.5$--$2$, in line with the threshold identified by Japkowicz and Stephen~\cite{japkowicz2002class}.
It is not merely an artefact of a smaller minority sample size, because the random removal procedure introduces no distributional bias into the degradation; the only systematic variable that changes is the relative count.
This directly addresses the long-standing open question of Weiss and Provost~\cite{weiss2003effect}: even when the minority distribution is held constant, imbalance alone degrades performance.

\textbf{However, the magnitude of the effect depends strongly on the classifier.}
Ensemble methods (RF)~\cite{breiman2001random, galar2012review} are substantially more robust than linear or distance-based methods (LR, KNN), suggesting that the ``imbalance problem'' is partly a problem of inappropriate algorithm choice~\cite{ling2004exploration}.
A well-tuned ensemble on a moderately imbalanced dataset ($\IR < 4$) may outperform SMOTE-augmented logistic regression, consistent with the findings of Galar et al.~\cite{galar2012review}.


% ============================================================================
\section{RQ5: Imbalance vs.\ Information Loss}
\label{sec:disc:rq5}

The random removal design was specifically intended to disentangle imbalance from information loss.
Because each instance is equally likely to be removed, the removal procedure introduces no distributional bias, so any performance loss must be attributable to the reduced \emph{count} (and thus the increased $\IR$) rather than to a systematic removal policy --- a confound identified as problematic by Japkowicz and Stephen~\cite{japkowicz2002class}.

Our results show that even with unbiased random removal, performance degrades monotonically with removal fraction.
This confirms that imbalance \emph{per se} --- not merely distributional shift --- is a genuine cause of degraded performance.
This extends the work of Weiss and Provost~\cite{weiss2003effect}, who demonstrated this effect for tree-based classifiers; our study shows it holds across five classifier families and six metrics.

Nonetheless, the recovery experiment reveals that real minority data cannot be fully replaced by synthetic data~\cite{blagus2013smote}.
When the starting minority is small (20\% remaining), all oversamplers recover to approximately 90--97\% of the balanced baseline, leaving a residual gap attributable to the irreducible information loss from having fewer real samples to learn from.
This is consistent with the theoretical bound of Blagus and Lusa~\cite{blagus2013smote}: SMOTE-synthesised distributions are a degraded approximation of the true minority distribution, and this approximation error propagates to the classifier's decision boundary.

\paragraph{Information-loss interpretation.}
When 80\% of real minority instances are removed, the surviving 20\% carry only a fraction of the original distributional information.
Even the best oversampler can only produce synthetic distributions whose fidelity to the original minority distribution is bounded by the information present in the surviving real instances.
This bound cannot be overcome by any synthesis method --- it is an information-theoretic floor, not an algorithmic limitation.
Part~V thus provides not only an empirical finding but a principled explanation for why the residual gap exists.


% ============================================================================
\section{Practical Implications}
\label{sec:disc:practical}

\begin{enumerate}[leftmargin=*, itemsep=4pt]
  \item \textbf{Oversampling is effective but not a panacea.}
  Users should expect recovery to within approximately 90--97\% of balanced performance, not full recovery~\cite{blagus2013smote, fernandez2018learning}.
  Collecting more real minority data remains preferable to synthetic augmentation where feasible.

  \item \textbf{Classifier choice matters as much as oversampling method.}
  Using an RF or ensemble classifier on the raw imbalanced data may be more practical than applying SMOTE to a weaker classifier, consistent with the recommendations of Galar et al.~\cite{galar2012review}.

  \item \textbf{GVM-CO and K-Means SMOTE offer the best recovery ceiling.}
  For applications where full recovery is important (e.g.\ medical diagnosis~\cite{he2009learning}), these methods are recommended over standard SMOTE.
  GVM-CO's advantage is mechanistically grounded: its gravity-biased Von Mises sampling (Chapter~\ref{ch:proposed_methods}) and BC-guided seed selection (Chapter~\ref{ch:seed_selection}) together preserve minority distributional fidelity more effectively than linear interpolation.

  \item \textbf{Leakage is a real and quantifiable concern.}
  Na\"{i}ve cross-validation that applies oversampling before splitting produces optimistically biased AUC estimates~\cite{kaufman2012leakage, cawley2010over}.
  The leakage-safe protocol described in Section~\ref{sec:setup:leakage} should be considered standard practice in oversampling research~\cite{johnson2019survey}.
  The quantified effect (2--5\% AUC inflation) is large enough to reverse the relative ranking of oversampling methods, making protocol standardization a prerequisite for reproducible benchmarking.
\end{enumerate}


% ============================================================================
\section{Limitations}
\label{sec:disc:limitations}

\begin{itemize}[leftmargin=*]
  \item The uniform random removal protocol is inherently stochastic: different random seeds produce different removal sequences.
  We use a fixed seed (\texttt{random\_state = 42}) for reproducibility, but the results may vary slightly under different seeds; the qualitative findings are expected to be robust~\cite{weiss2003effect}.

  \item Real KEEL/UCI files were not available for all datasets; synthetic analogues with matching $(n, d)$ were used~\cite{blagus2013smote}.
  The qualitative findings are expected to generalise~\cite{fernandez2018learning}, but exact numerical results may differ on the original files.

  \item The recovery experiment starts from $\IR \approx 4$.
  Behavior at more extreme ratios ($\IR > 10$~\cite{he2009learning}) is not directly characterized here; the threshold effect and recovery gap may be more pronounced at higher imbalance.

  \item Only binary classification is studied.
  Multi-class imbalance introduces additional complexity --- including competing minority-to-minority interactions --- not addressed here~\cite{fernandez2018learning, haixiang2017learning}.
\end{itemize}
